\begin{description}
\item[(a)] Let $v_p$ be the closest approach velocity of the star. From
  conservation of mechanical energy:
  \[ -\frac{G M_T m}{r_p} + \frac{1}{2} m v_p^2 = \frac{1}{2} m v^2 \]
  \hfill(2.0)

  where $M_T = 2M$ is the total mass of the binary. From conservation of
  angular momentum, we obtain $v_p$:
  \[ v_p = \frac{b}{r_p}\,v \]
  \hfill(2.0)

  setting $r_p \approx a/2$ and inserting $v_p$ into the first equation, we
  solve for $b$:
  \begin{align*}
    b &= a\sqrt{\frac{1}{4} + \frac{2GM}{a v^2}} \\
    b &\approx \frac{\sqrt{2GMa}}{v}
  \end{align*}
  \hfill(1.0)

  Where we have used $v^2 \ll \frac{GM}{a}$ in the last step.

  \textbf{Note:} The student can also solve it by exploring the geometry of
  the star's orbit, which is hyperbolic, by writing
  $r_p = a_{orbit}(e - 1)$ and finding $a_{orbit}$ and $e$ by writing out
  equations for the angular momentum and mechanical energy.

\item[(b)] The speed of the black hole in circular orbit can be found by
  equating the gravitational force to the centripetal net force:
  \[ \frac{GM^2}{a^2} = \frac{M V^2}{\frac{a}{2}}
     \qquad \therefore\ V = \sqrt{\frac{GM}{2a}} \]
  \hfill(1.0)

  For the rest of the problem, we propose two different solutions:

  \textbf{Solution I:}
  Additionally, let $v_f$ be the ejection speed of the star, and $V'$ the
  final speed of the black hole. To find $v_f$, we use conservation of linear
  momentum and conservation of mechanical energy for the two bodies:
  \[ M V = M V' + m v_f \]
  \hfill(1.5)
  \[ \frac{1}{2} M V^2 = \frac{1}{2} m v_f^2 + \frac{1}{2} M V'^2 \]
  \hfill(1.5)
  \begin{align*}
    \frac{1}{2} M V^2
      &= \frac{1}{2} m v_f^2
       + \frac{1}{2} M \left( V - \frac{m}{M} v_f \right)^2 \\
    0 &= v_f^2 - 2 V v_f + \frac{m}{M} v_f^2
  \end{align*}
  Since $m \ll M$ \hfill(1.0)\\
  we neglect the last term, yielding:
  \[ v_f = 2V = \sqrt{\frac{2GM}{a}} \]
  \hfill(1.0)

  Because we are considering two moments where the distance between the star
  and the black hole is very large, we disregard potential energy terms.

  \textbf{Solution II:}
  In the reference frame of the component, the star approaches the component
  with a velocity $-\vec{V}$. \hfill(1.5)

  And, since mechanical energy and linear momentum are conserved, it can be
  shown that the relative speed between the components is the same before and
  after the interaction --- that is, restitution coefficient equal to unity.
  Furthermore, since $M \gg m$, the velocity of the component can be
  considered unchanged (notice this is not valid in Solution I, as there we
  deal with second order terms arising from energy considerations) in the
  interaction. Therefore, after the interaction, the star is ejected back with
  a velocity $\vec{v}\,'_f = \vec{V}$ in the opposite direction. \hfill(1.5)

  To return to the CM frame, we write $\vec{v}_f = \vec{v}\,'_f + \vec{V}$,
  such that
  \[ v_f = 2V \]
  \hfill(2.0)

  As we have previously found.

\item[(c)] First, we know the kinetic energy acquired by a star during an
  encounter:
  \[ \Delta K_* = \frac{1}{2} m v_f^2 - \frac{1}{2} m v^2
     \approx \frac{GMm}{a} \]
  \hfill(1.0)

  Where we used that $v^2$ is negligible with respect to $v_f^2$, since
  $v_f^2 = \frac{2GM}{a}$, and $v^2 \ll \frac{GM}{a}$. This is also the
  energy lost by the binary during an encounter.

  In order to find the rate of energy extraction, we must first obtain the
  rate of encounters, which can be estimated as follows: imagine the binary
  travelling with velocity $v$ in the reference frame of a far away star of
  impact parameter $b$. During a small interval $dt$, all the stars included
  inside the annular cylinder of length $L = v\,dt$ bounded by radii of $b$
  and $b + db$ will interact with the star. Call this number of stars $dN$.
  It is given by:
  \begin{align*}
    dN &= n \cdot V_{\mathrm{cylinder}} \\
       &= n \cdot A_{\mathrm{annulus}} \cdot L \\
       &= n \cdot \pi \cdot \left[ (b + db)^2 - b^2 \right] \cdot v \cdot dt \\
       &\approx 2\pi n\, b\, v\, db\, dt
  \end{align*}
  \hfill(4.0)

  Defining $d\varphi \equiv \frac{dN}{dt}$ to be the rate of encounters with
  stars of impact parameter between $b$ and $b + db$ and using $b$ from part
  (a),
  \begin{align*}
    d\varphi &= 2\pi n\, b\, v\, db \\
             &= 2\pi n \cdot \frac{\sqrt{2GMa}}{v} \cdot v \cdot db \\
    d\varphi &= 2\pi n \sqrt{2GMa} \cdot db
  \end{align*}
  \hfill(1.0)

  To find the total rate of encounters ($\varphi = \sum d\varphi$), we sum
  over all impact parameters in the interval
  $\frac{1}{2} b_0 \le b \le \frac{3}{2} b_0$. Notice
  $\sum db = \frac{3}{2} b_0 - \frac{1}{2} b_0 = b_0$, so that:
  \begin{align*}
    \varphi &= 2\pi n \sqrt{2GMa} \cdot \sum db \\
            &= 2\pi n \sqrt{2GMa} \cdot \frac{\sqrt{2GMa}}{v_0} \\
            &= \frac{4\pi n G M a}{v_0}
  \end{align*}
  \hfill(3.0)

  Since each encounter takes away $\Delta K_*$ and there are $\varphi$
  encounters per unit time, the rate of change of energy for the binary is:
  \[ \frac{dE_{\mathrm{bin}}}{dt} = (-\Delta K_*) \cdot \varphi
     = -\frac{4\pi n G^2 M^2 m}{v_0} \]
  \hfill(2.0)

  A change in the binary total energy can also be written as:
  \[ \frac{dE_{\mathrm{bin}}}{dt}
     = \frac{d}{dt}\left( -\frac{G M^2}{2a} \right)
     = -\frac{G M^2}{2} \frac{d}{dt}\left( \frac{1}{a} \right) \]
  \hfill(2.0)

  Therefore, equating the two:
  \[ \frac{d}{dt}\left( \frac{1}{a} \right) = \frac{8\pi G n m}{v_0} \]
  \[ \frac{d}{dt}\left( \frac{1}{a} \right) = 8\pi\,\frac{G\rho}{v_0} \]
  \hfill(1.0)

  Thus, $H = 8\pi$.
\end{description}
